\documentclass[10pt]{article}

\usepackage[utf8]{inputenc}

\usepackage[T1]{fontenc}

\usepackage{amsmath}

\usepackage{amsfonts}

\usepackage{amssymb}

\usepackage[version=4]{mhchem}

\usepackage{stmaryrd}

\usepackage{bbold}



\begin{document}

ФУНКЦИОНАЛЬНАЯ ГИПОТЕЗА - см. Сокращенное описание состояния.



ФУНКЦИОНАЛЬНАЯ ПРОИЗВОДНАЯ - см. Вариационное исчисление.



ФУНКЦИОНАЛЬНЫЙ ИНТЕГРАЛ, инте грал по траекториям, континуальный интеграл,- интеграл, для которого область определения - то или иное функциональное пространство. Ф. и. можно определить как интеграл Лебега от функционала, заданного на пространстве функций, по нек-рой мере в этом пространстве. Такие меры можно получить распространением кваз и мер (то есть аддитивных функций множеств, определенных на алгебре всех цилиндрич. подмножеств функционального пространства). Ф. и. определяют также как предел по мере Лебега в $\mathbb{R}^{n}$ при $n \rightarrow \infty$. При этом $\mathbb{R}^{n}$ служит приближением для функционального пространства, а приближением функционала является функция $n$ переменных.



Многие из Ф. и., встречающихся в теоретич. физике, пока не имеют строгого определения, и операции над ними часто производятся на формальном уровне без строгого обоснования. С Ф. и. в квантовой теории поля и квантовой статистич. физике связаны проблемы общей теории меры и интегрирования по бесконечномерному пространству.



Ф. и. в к вантовой механ и ке. Матричный элемент оператора эволюции $\exp (-i \hat{H} t), t \in \mathbb{R}^{1}$, квантовомеханич. системы можно записать в виде Ф. и., определив его как предел $n$-кратного интеграла при $n \rightarrow \infty$.



Пусть $H(p, q)$ - функция Гамильтона одномерной механич. системы ( $q$ - координата, $p-$ импульс),



\[

S=\int_{0}^{t}(\dot{p}(\tau) q(\tau)-H(p(\tau), q(\tau)) d \tau

\]



\begin{itemize}

  \item соответствующий функционал действия $(\dot{q}(\tau)=d q(\tau) / d \tau)$. Интервал $[0, t]$ делится на $N$ подинтервалов точками деления $\tau_{1}, \ldots, \tau_{N-1}$ и рассматриваются функции $p(\tau)$, постоянные на подинтервалах, и функции $q(\tau)$, линейные на подинтервалах. Такая траектория определяется значениями $p_{1}$, $\ldots, p_{N}$ функции $p(\tau)$ на подинтервалах и значениями $q=q_{0}$, $q_{1}, \ldots, q_{N-1}, q_{N}=q_{1}$ функции $q(\tau)$ на концах подинтервалов. Фиксируя $q(0)=q_{0}=q, q(t)=q_{N}=q^{\prime}$, вводят интеграл

\end{itemize}



\[

\begin{gathered}

J_{N}\left(q, q^{\prime}, t\right)=(2 \pi)^{-N} \int \prod_{i=1}^{N} d p_{i} \times \\

\times \prod_{j=1}^{N-1} d q_{j} \exp i S\left(p_{1}, \ldots, p_{N}, q_{0}, q_{1}, \ldots, q_{N}\right) .

\end{gathered}

\]



В пределе при $N \rightarrow \infty$ этот интеграл дает матричный элемент оператора инволюции, а именно,



\[

\lim _{N \rightarrow \infty} \int d q d q^{\prime} f(q) J_{N}\left(q, q^{\prime}, t\right) \bar{g}\left(q^{\prime}\right)=\left(e^{-i \hat{H} t} f, g\right) .

\]



Если $H(p, q)$ зависит только от $p$ или только от $q$, допредельное выражение для $J_{N}$ не зависит от $N$ и совпадает со своим пределом. Если $H=p^{2} / 2 m+V(q)$, то можно доказать существование предела, используя формулу Троттера-Като. Для функции Гамильтона $H(p, q)$ общего вида предел $J_{N}$ зависит от способа аппроксимации траектории $(p(\tau), q(\tau))$. Это связано с проблемой построения оператора $\hat{H}$ по функции Гамильтона $H(p, q)$ заменой $p, q$ операторами $\hat{p}, \hat{q}$ с перестановочными соотношениями $\hat{q} \hat{p}-\hat{p} \hat{q}=i$ (проблема упорядочения, не имеющая, как известно, естественного решения).



Рассмотренная конструкция Ф. и. непосредственно переносится на гамильтонову систему с любым конечным числом степеней свободы, где $q=\left(q^{1}, \ldots, q^{n}\right)$ - точка многообразия $M, p=\left(p_{1}, \ldots, p_{n}\right)-$ элемент кокасательного пространства.



Если заменить в (1) $H(p(\tau), q(\tau))$ на $-i H(p(\tau), q(\tau))$, то выражение $J_{N}$ в (2) в пределе $N \rightarrow \infty$ дает матричные элементы для полугруппы операторов $\exp (-\hat{H} t), t>0$. В част- ном случае $H=D p^{2}+V(q)$ получается $\Phi$. и. для уравнения теплопроводности - исторически первого объекта, для к-рого возникло понятие Ф. и. (см. [1]).



Ф.и. в квантовой статистической физике. Для системы бозе-частиц в кубич. объеме $V=L^{D}$ с периодич. граничными условиями Ф. и.- интеграл с весом $\exp S$ по пространству комплексных функций $\Psi(x, \tau)$ (где $x \in V$, $\tau \in[0, \beta], \beta^{-1}=T-$ абсолютная температура), периодических по $\tau: \Psi(x, \beta)=\Psi(x, 0)$ (см. [7], [8]). Функционал действия $S$ имеет вид:



\[

S=\int_{0}^{\beta} d \tau \int d x \bar{\Psi}(x, \tau) \partial \tau \Psi(x, \tau)-\int_{0}^{\beta} H^{\prime}(\tau) d \tau,

\]



где



\[

\begin{aligned}

& H^{\prime}(\tau)=\int d x\left((2 m)^{-1} \nabla \Psi^{\prime}(x, \tau) \nabla \Psi(x, \tau)-\lambda \bar{\Psi}(x, \tau) \Psi(x, \tau)\right)+ \\

& \quad+\frac{1}{2} \int d x d y u(x-y) \bar{\Psi}(x, \tau) \bar{\Psi}(y, \tau) \Psi(y, \tau) \Psi(x, \tau)

\end{aligned}

\]



\begin{itemize}

  \item гамильтониан, в к-ром $\lambda$ - химич. потенциал, $u(x-y)-$ парный потенциал взаимодействия между бозе-частицами. Разложив $\Psi$ в ряд Фурье:

\end{itemize}



\[

\Psi(x, \tau)=(\beta V)^{-1 / 2} \sum_{p} e^{i p x} a(p),

\]



где $p=(k, \omega), p x=\omega \tau+k x, \omega=\omega_{n}=2 \pi n / \beta-$ частота Мацуба$\mathrm{pa}, k_{i}=2 \pi n_{i} / L$, можно затем определить $\Phi$. и. как интеграл от нек-рого функционала $F(\Psi, \bar{\Psi})$ по мере



\[

D_{\mu}=\prod_{p} d a^{*}(p) d a(p)

\]



с весом $\exp S$, где $S$ в терминах $a(p), a^{*}(p)$ имеет вид:



\[

\begin{aligned}

& S=\sum_{p}\left(i \omega-k^{2} / 2 m+\lambda\right) a^{*}(p) a(p)- \\

& -(4 \beta V)^{-1} \sum_{p_{1}+p_{2}=p_{3}+p_{4}}\left(v\left(k_{1}-k_{3}\right)+\right. \\

& \left.+v\left(k_{1}-k_{4}\right)\right) a^{*}\left(p_{1}\right) a^{*}\left(p_{2}\right) a\left(p_{4}\right) a\left(p_{3}\right) .

\end{aligned}

\]



При этом $v(k)=\int d x e^{i k x} u(x)-$ потенциал «в $k$-представлении». Такой $\Phi$. и. есть предел конечномерного интеграла по переменным $a(p), a^{*}(p)$ с $|k|<K_{0},|\omega|<\Omega_{0}$ при $K_{0}, \Omega_{0} \rightarrow \infty$. Функции Грина определяются как (формальные) частные Ф. и.:



\[

G=-\left\langle\Psi(x, \tau) \bar{\Psi}\left(x^{\prime}, \tau^{\prime}\right)\right\rangle=-\int D \mu e^{S} \Psi(x, \tau) \bar{\Psi}\left(x^{\prime}, \tau^{\prime}\right) / \int D_{\mu} e^{S} .

\]



Теория возмущений для вычисления функций Грина естественно возникает, если считать невозмущенным действием квадратичную форму по $a(p), a^{*}(p)$ в $S(3)$, а форму четвертой степени считать возмущением. Каждому члену теории возмущений можно сопоставить диаграмму, состоящую из элементов - линий и вершин. Линии соответствует невозмущенная функция Грина $G_{0}(p)=\left(i \omega-k^{2} / 2 m+\lambda\right)^{-1}$, вершине - симметризованный потенциал $v\left(k_{1}-k_{3}\right)+$ $+v\left(k_{1}-k_{4}\right)$. Выражение для диаграммы, дающей вклад в функцию Грина, получают, просуммировав произведение выражений для ее линий и вершин по независимым 4-импульсам $p=(k, \omega)$ внутренних линий и умножив результат на $(-1 / \beta v)^{l-v-1} r^{-1}$, где $l-$ число линий, $v-$ число вершин, $r$ - порядок группы симметрии диаграммы.



Заменив в описанной конструкции Ф. и. комплексные числа $a(p), a^{*}(p)$ в (3) антикоммутируюшими друг с другом образующими грассмановой алгебры, получают Ф. и. для системы ферми-частиц в объеме $V$ при конечной температуpe $T$, определяемый как предел конечномерного интеграла по антикоммутирующим переменным $a(p), a^{*}(p) \mathrm{c}|k|<K_{0}$, $|\omega|<\Omega_{0}$ (интеграл Березина) при $K_{0}, \Omega_{0} \rightarrow \infty$. В выражении для линии $\left(i \omega-k^{2} / 2 m+\lambda\right)^{-1}$ диаграммной теории возмуще-





\end{document}